\subsection{Phiên bản sản phẩm tối thiểu khả thi, gọi là MVP}

MVP của 1Stream được xác định là phiên bản tối thiểu có khả năng trình diễn và
kiểm chứng luồng giá trị cốt lõi của sản phẩm: từ dữ liệu khóa học đã được chuẩn
hóa, hệ thống tạo nội dung video AI, hỗ trợ vận hành một phiên livestream thử
nghiệm, trả lời các câu hỏi tuyển sinh phổ biến và chuyển các trường hợp ngoài
phạm vi cho nhân viên phụ trách.

MVP không nhằm triển khai ngay một nền tảng livestream AI đa kênh hoàn chỉnh
hoặc tự động hóa toàn bộ quy trình. Mục tiêu của phiên bản này là kiểm chứng ba
giả định chính:

\begin{itemize}
    \item trung tâm ngoại ngữ chấp nhận sử dụng video AI cho các nội dung tuyển
    sinh có tính lặp lại;

    \item AI Agent có thể hỗ trợ trả lời một tập câu hỏi phổ biến dựa trên dữ
    liệu đã được khách hàng xác nhận;

    \item quy trình kết hợp giữa video AI, hỗ trợ tương tác và chuyển lead có
    thể giúp giảm khối lượng công việc lặp lại trong một phiên livestream.
\end{itemize}

Trong giai đoạn MVP, 1Stream được triển khai theo mô hình
\textit{service-enabled software}. Một số bước như kiểm tra dữ liệu, duyệt nội
dung, cấu hình phiên phát và giám sát sự cố vẫn được đội ngũ 1Stream thực hiện
thủ công nhằm kiểm soát chất lượng và thu thập dữ liệu trước khi tự động hóa ở
các phiên bản tiếp theo.

\subsubsection*{Phạm vi khách hàng thử nghiệm}

MVP ưu tiên thử nghiệm với một nhóm nhỏ trung tâm tiếng Anh, tiếng Hàn hoặc đơn
vị bán khóa học ngôn ngữ có hoạt động trên Facebook hoặc TikTok, có dữ liệu
khóa học tương đối đầy đủ và có người phụ trách duyệt nội dung, vận hành
livestream hoặc tiếp nhận lead.

Dữ liệu đầu vào tối thiểu gồm:

\begin{itemize}
    \item tên và mô tả khóa học;
    \item học phí, lịch khai giảng và ưu đãi;
    \item lộ trình học và trình độ đầu vào;
    \item thông tin đăng ký học thử hoặc tư vấn;
    \item các câu hỏi tuyển sinh thường gặp;
    \item hình ảnh, video hoặc tài sản truyền thông có quyền sử dụng.
\end{itemize}

\subsubsection*{Các thành phần cốt lõi của MVP}

\begin{enumerate}
    \item \textbf{Tiếp nhận và chuẩn hóa dữ liệu khóa học:}
    cung cấp biểu mẫu nhập học phí, lịch khai giảng, lộ trình, ưu đãi, CTA và
    FAQ. Các thông tin thiếu hoặc mâu thuẫn được đánh dấu để khách hàng xác nhận
    trước khi sử dụng.

    \item \textbf{Tạo nội dung video AI:}
    tạo bản nháp kịch bản từ dữ liệu khóa học, kết hợp với một giọng đọc hoặc
    avatar có sẵn, hình ảnh, video nền và phụ đề để tạo một video livestream
    mẫu có khả năng phát lặp.

    \item \textbf{Quy trình duyệt nội dung:}
    khách hàng xem trước, phản hồi và xác nhận phiên bản nội dung cuối cùng.
    Chỉ video và dữ liệu đã được duyệt mới được đưa vào phiên thử nghiệm.

    \item \textbf{Quản lý một phiên livestream thử nghiệm:}
    hiển thị video AI, trạng thái phiên và luồng bình luận trong môi trường demo
    có kiểm soát hoặc trên một kênh thử nghiệm đủ điều kiện.

    \item \textbf{AI Agent hỗ trợ FAQ:}
    nhận diện và trả lời một số nhóm câu hỏi cơ bản như học phí, lịch khai
    giảng, trình độ đầu vào, ưu đãi và đăng ký học thử dựa trên cơ sở tri thức
    đã được xác nhận.

    \item \textbf{Chuyển người thật và ghi nhận lead:}
    đánh dấu các câu hỏi ngoài cơ sở tri thức, câu hỏi chuyên môn hoặc nội dung
    nhạy cảm để chuyển cho nhân viên. Người xem có nhu cầu tư vấn hoặc đăng ký
    học thử được ghi nhận dưới dạng lead cơ bản.

    \item \textbf{Báo cáo sau phiên:}
    tổng hợp thời lượng phát, số lượt bình luận, số phản hồi AI, số câu hỏi cần
    chuyển người thật, số lead được ghi nhận và các lỗi phát sinh.
\end{enumerate}

\subsubsection*{Đầu ra của MVP}

Đầu ra của MVP được xác định theo các thành phần có thể trình diễn, kiểm tra và
đo lường trực tiếp, bao gồm:

\begin{itemize}
    \item một bộ dữ liệu khóa học mẫu đã được chuẩn hóa;

    \item một kịch bản và video livestream AI hoàn chỉnh sử dụng avatar hoặc
    giọng đọc có sẵn;

    \item một phiên livestream thử nghiệm trong môi trường demo có kiểm soát
    hoặc trên một kênh đủ điều kiện;

    \item một cơ sở tri thức chứa các câu hỏi tuyển sinh phổ biến;

    \item một AI Agent có khả năng phản hồi tập FAQ cơ bản;

    \item một cơ chế chuyển câu hỏi ngoài phạm vi cho người thật;

    \item một chức năng ghi nhận lead cơ bản;

    \item một báo cáo sau phiên gồm dữ liệu vận hành, tương tác và chuyển giao;

    \item một bộ kết quả kiểm thử gồm thời gian tạo video, tỷ lệ phản hồi đúng,
    tỷ lệ phản hồi cần chỉnh sửa, tỷ lệ phải chuyển người thật và các lỗi còn
    tồn tại.
\end{itemize}

MVP được xem là hoàn thành khi nhóm có thể trình diễn liên tục luồng sau:

\begin{align*}
&\text{Dữ liệu khóa học} \rightarrow \text{Chuẩn hóa} \rightarrow \text{Video AI} \\
&\rightarrow \text{Duyệt nội dung} \rightarrow \text{Phiên livestream thử nghiệm} \\
&\rightarrow \text{AI xử lý FAQ} \rightarrow \text{Chuyển lead} \rightarrow \text{Báo cáo}
\end{align*}

\subsubsection*{Các chức năng chưa thuộc phạm vi MVP}

Để tránh mở rộng phạm vi quá sớm, MVP chưa bao gồm:

\begin{itemize}
    \item phát đồng thời trên nhiều nền tảng ở quy mô lớn;
    \item cam kết vận hành liên tục 24/7;
    \item avatar tùy chỉnh phức tạp hoặc sao chép hoàn toàn hình ảnh người thật;
    \item hỗ trợ nhiều ngôn ngữ trong cùng một phiên;
    \item tự động tư vấn chuyên môn học thuật;
    \item tự động chốt thanh toán hoặc xác nhận đăng ký khóa học;
    \item tích hợp sâu với CRM, ERP hoặc hệ thống quản lý đào tạo;
    \item phân quyền doanh nghiệp, SLA và các yêu cầu Enterprise;
    \item nền tảng SaaS tự phục vụ hoàn toàn không có đội ngũ hỗ trợ.
\end{itemize}

Các chỉ số dùng để đánh giá MVP gồm thời gian chuẩn hóa dữ liệu, thời gian tạo
video, mức độ chấp nhận video AI, tỷ lệ phản hồi FAQ chính xác, tỷ lệ câu hỏi
cần chuyển người thật, độ ổn định của phiên thử nghiệm và phản hồi của khách
hàng về khả năng giảm công việc lặp lại. Mức sẵn sàng tham gia paid pilot hoặc
mua gói được sử dụng để kiểm chứng thị trường và quyết định chuyển sang giai
đoạn phát triển tiếp theo, không phải là điều kiện kỹ thuật để xác định MVP đã
hoàn thành.